% Source: https://artofproblemsolving.com/community/c3h1715535_need_help
% Stable identity: aops-topic-1715535-q01
Let $AD=BC=x$. In a right triangle each leg is the geometric mean of the hypotenuse and its projection on the hypotenuse (triangles $CBD$ and $ABC$ are similar, as are $ACD$ and $ABC$):
\[
BC^2=BD\cdot AB,\qquad AC^2=AD\cdot AB .
\]

\begin{asy}
usepackage("amsmath");
size(6cm);
defaultpen(fontsize(10pt));
void ra(pair P,pair Q,pair R,real s){pair u=unit(Q-P),v=unit(R-P);draw(P+s*u--P+s*u+s*v--P+s*v,linewidth(0.5));}
real x=(3*sqrt(5)-3)/2; real h=sqrt(x*(3-x));
pair A=(0,0),B=(3,0),D=(x,0),C=(x,h);
draw(A--B--C--cycle,linewidth(0.9));draw(C--D,linewidth(0.7));
ra(D,B,C,0.12);ra(C,A,B,0.14);
pair m1=(A+D)/2,m2=(B+C)/2,n2=unit(rotate(90)*(C-B));
draw(m1+(0,-0.07)--m1+(0,0.07));draw(m2-0.07*n2--m2+0.07*n2);
label("$A$",A,SW);label("$B$",B,SE);label("$C$",C,N);label("$D$",D,S);
label("$3$",(1.5,-0.12),S);
label("\mbox{}",(min(currentpicture,user=true)+max(currentpicture,user=true))/2,black); // keeps the text after the figure black when LaTeX typesets the labels inline
\end{asy}

With $BD=3-x$ the first relation gives $x^2=3(3-x)$, i.e.\ $x^2+3x-9=0$, so $x=\dfrac{3\sqrt5-3}{2}$ (the positive root). Then
\[
AC^2=3x=\dfrac{9\left(\sqrt5-1\right)}{2},\qquad AC=\boxed{3\sqrt{\dfrac{\sqrt5-1}{2}}}\ \text{cm}\approx2.36\text{ cm}.
\]
Check: $AC^2+BC^2=3x+x^2=3x+(9-3x)=9=AB^2$.

(The post also gives $CD=\sqrt2$, but this contradicts the other data: from $AD=BC$ and $AB=3$ we get $CD^2=AD\cdot DB=x(3-x)=9\sqrt5-18\approx2.12$, not $2$. With $CD=\sqrt2$ and $AB=3$ instead, $\{AD,DB\}=\{1,2\}$ and $BC=\sqrt6$ or $\sqrt3$, never equal to $AD$.)
